% FRAGMENT A12 (drafter A: class model and main results)
% REPLACES: main.tex lines 1254–1397 ("\section{Appendix A}\label{paper-appendix}" to "… conditional-probability argument. \(\square\)")
%   This covers the Appendix A opening and Appendix A.1 (proof of Proposition 1,
%   Propositions A.1 and A.2).
% LABELS KEPT: paper-appendix, pa-results, proof-of-proposition-1, eq:paper-a-1,
%   eq:paper-a-2, posterior-bounds, result-proposition-a-1,
%   price-sufficiency-and-residual-profits, result-proposition-a-2, eq:paper-a-3.
% LABELS REMOVED: none.  LABELS ADDED: none.
% TAGS: (A.1)–(A.3) keep their numbers and their objects: (A.1) the point-mass uniform
%   integrals, (A.2) the indicator integrals (now with the class kernel), (A.3) the price
%   increment. All other new displays are unnumbered.
% NEW SYMBOLS: \xi(R) (expected difference in proceeds between the classes at incumbent
%   value R; local to the proof).
% NEW CITATIONS: none.  NUM PLACEHOLDERS: none.
% SOURCES: Q1_N_main.md Claims 1–4 (N1)–(N9), 6, 7; SYNTHESIS.md A.1 rows 1–3 (including
%   the correction that the F-free split also holds when F(p)=0).
% ONLINE APPENDIX POINTERS kept as in the manuscript (A.1, A.2, A.3). Those Online Appendix
%   sections are written for the point-mass case; drafter B's scope note B05 records which
%   arguments carry over to classes.

\section{Appendix A}\label{paper-appendix}

This appendix proves every result in the paper. Each statement
carries its status once: analytical, computer-assisted, numerical
diagnostic, or open. Statements use the class model of
\hyperref[sec-model]{Section 2}. A statement that needs the
point-mass benchmark says so. The Online Appendix holds the
probability-space construction, regular conditional distributions,
null-set invariance under deviations, convolution regularity, the
Stieltjes argument, interval endpoint ordering, the complete signal
laws, and the numerical definitions. Equations \eqref{eq:paper-1} to
\eqref{eq:paper-14} are in the main text.

\subsection{A.1 Acquisition payoffs and Bayesian
pricing}\label{pa-results}

\subsubsection{\texorpdfstring{Proof of
\hyperref[result-proposition-1]{Proposition
1}}{Proof of Proposition 1}}\label{proof-of-proposition-1}

Without entry the target receives \(p\mathbf1\{R\ge p\}\). A prepared
challenger with \(\theta\ge p\) wins at \(\theta>R\) and pays
\(\max\{p,R\}\). A loss makes the incumbent pay \(\theta\). The target
receives \(\max\{p,\min(R,\theta)\}\) and the challenger keeps
\((\theta-\max\{p,R\})_+\). At \(\theta<p\) the challenger does not
bid. The target again receives \(p\mathbf1\{R\ge p\}\) and the
challenger keeps nothing. Truthful bidding is weakly dominant at
every competing bid. No payoff here rests on a shading conjecture or
a belief. In the point-mass benchmark with the uniform incumbent,

\[
\begin{aligned}
t_H&=\frac1r\left[\int_0^p p\,du+\int_p^r u\,du\right],\\
t_L&=\frac1r\left[\int_0^p p\,du+\int_p^\ell u\,du+\int_\ell^r\ell\,du\right],\\
g_L&=\frac1r\left[\int_0^p(\ell-p)\,du+\int_p^\ell(\ell-u)\,du\right],
\end{aligned}
\tag{A.1}\label{eq:paper-a-1}
\]

which evaluate to \eqref{eq:paper-4} and differentiate to
\eqref{eq:paper-5}.

For general classes and a continuous \(F\), write the payments as
integrals of indicators. For \(\theta\ge p\),

\[
\max\{p,\min(R,\theta)\}=p+\int_p^\theta\mathbf1\{R>u\}\,du,
\qquad
(\theta-\max\{p,R\})_+=\int_p^\theta\mathbf1\{R<u\}\,du.
\tag{A.2}\label{eq:paper-a-2}
\]

The two indicators sum to one for almost every \(u\), which gives
\(t(\theta)+g(\theta)=\theta\). Tonelli's theorem applied to the
nonnegative integrands gives
\(t(\theta)=p+\int_p^\theta[1-F(u)]\,du\) and
\(g(\theta)=\int_p^\theta F(u)\,du\) for \(\theta\ge p\), while
\(t(\theta)=t_0=p[1-F(p)]\) and \(g(\theta)=0\) for \(\theta<p\). At
\(\theta=p\), \(t\) jumps by \(pF(p)\): when the incumbent fails the
reserve, an admissible challenger buys at \(p\) and an inadmissible one
leaves the target unsold. Hence, for every \(\theta\),

\[
t(\theta)=t_0+pF(p)\mathbf1\{\theta\ge p\}
+\int_p^\infty\mathbf1\{u\le\theta\}[1-F(u)]\,du,\qquad
g(\theta)=\int_p^\infty\mathbf1\{u\le\theta\}F(u)\,du .
\]

Take the class expectation and apply Tonelli's theorem again, with
\(\mathbb E_j\mathbf1\{\theta\ge p\}=1-G_j(p^-)\) and
\(\Pr_j(\theta\ge u)=1-G_j(u)\) for almost every \(u\). This gives

\[
t_j=t_0+pF(p)[1-G_j(p^-)]+\int_p^\infty[1-G_j(u)][1-F(u)]\,du
\]

and \(g_j\) in \eqref{eq:paper-6}; subtracting the classes gives
\(\Delta_T\) in \eqref{eq:paper-6}. The finite class means make every
integral finite. Since
\(\mathbb E_j[\max\{\theta,p\}]=p+\int_p^\infty[1-G_j(u)]\,du\), adding
the two lines of \eqref{eq:paper-6} gives the split in its general form,

\[
\Delta_T(F)+g_H(F)-g_L(F)
=pF(p)D(p^-)+\mathbb E_H[\max\{\theta,p\}]-\mathbb E_L[\max\{\theta,p\}].
\]

The reserve term vanishes under equal admissibility, and also when the
incumbent always meets the reserve, \(F(p)=0\).

\emph{Part (a).} In \eqref{eq:paper-6} the weight \(1-G_j\) on \(F\) is
nonnegative, so the pointwise decrease from \(F\) to \(\tilde F\) lowers
\(g_j\), strictly exactly when
\(\int_p^\infty[1-G_j][F-\tilde F]\,du>0\). Gross profit at a fixed
belief is an average of \(g_H\) and \(g_L\) with fixed weights, so it
falls as well.

\emph{Part (b).} Subtracting \eqref{eq:paper-6} at the two
distributions,

\[
\Delta_T(\tilde F)-\Delta_T(F)
=\int_p^\infty[F(u)-\tilde F(u)]D(u)\,du-p[F(p)-\tilde F(p)]D(p^-).
\]

Under the restrictions of
\hyperref[sec-values-preparation-and-the-sale]{Section 2.2} the second
term is zero and the integrand is nonnegative, so the spread rises,
strictly exactly when the integral is positive.

\emph{Part (c).} Fix the incumbent's value \(R\). If \(R\ge p\), the
target receives \(\min\{R,\max\{\theta,p\}\}\) for every challenger
value: an admissible challenger and the incumbent both bid at least
\(p\), and an inadmissible challenger leaves the incumbent to pay \(p\).
If \(R<p\), the target receives \(p\mathbf1\{\theta\ge p\}\). Since
\(\mathbb E_j[\min\{R,\max\{\theta,p\}\}]=p+\int_p^R[1-G_j(u)]\,du\)
for \(R\ge p\), the expected difference in proceeds between the classes
at incumbent value \(R\) is

\[
\xi(R)=
\begin{cases}
pD(p^-), & R<p,\\[3pt]
\int_p^RD(u)\,du, & R\ge p,
\end{cases}
\]

and \(\Delta_T(F)=\mathbb E_F[\xi(R)]\). The function \(\xi\)
is bounded, constant below \(p\), continuous on \([p,\infty)\) with
derivative \(D\) almost everywhere, and moves from \(pD(p^-)\) to
\(\xi(p)=0\) at \(p\). It is nondecreasing if and only if
\(D\ge0\) almost everywhere above \(p\) and \(D(p^-)\le0\). If
\(\xi\) is nondecreasing, every first-order stochastic
strengthening raises \(\mathbb E_F[\xi(R)]\). If it is not,
continuity on each side of \(p\) gives points \(0<x_1<x_2\), both
different from \(p\), with \(\xi(x_1)>\xi(x_2)\). Let \(F\) be
uniform on \([x_1-\eta,x_1+\eta]\) and let \(\tilde F\) be the same law
shifted to \(x_2\). Then \(\tilde F\le F\), and as \(\eta\to0\),
\(\Delta_T(\tilde F)-\Delta_T(F)\to\xi(x_2)-\xi(x_1)<0\). If
the support must start at zero, mix the same small uniform component on
\([0,\bar r]\) into both laws; this scales the difference by a positive
factor and keeps its sign. Because distribution functions are
right-continuous, \(D\ge0\) almost everywhere above \(p\) is the same
as \(D(u)\ge0\) for every \(u>p\), and it implies \(D(p)\ge0\). When
neither class has an atom at \(p\), \(D(p^-)=D(p)\ge0\), so the
condition holds exactly when \(D(p^-)=0\), which with the ranking above
the reserve is the pair of restrictions of
\hyperref[sec-values-preparation-and-the-sale]{Section 2.2}.

For the point masses with \(p<\ell\), \(D(u)=\mathbf1\{\ell\le u<h\}\),
\(D(p^-)=0\), and \(\xi(R)=(R-\ell)_+\) for \(R<h\). Equation
\eqref{eq:paper-6} then gives \(\Delta_T=\int_\ell^h[1-F(u)]\,du\),
or \(\int_\ell^{\bar r}[1-F(u)]\,du\) at \(\bar r<h\), and
\(g_j=\int_p^\theta F(u)\,du\) at \(\theta=h\) and \(\theta=\ell\).
The kernel \(\xi\) belongs to the cash second-price rule with reserve.
\hyperref[sec-the-payment-rule]{Section 6.2} changes the rule and the
kernel with it. Online Appendix A.2 gives the measure formulation.
\(\square\)

\Needspace{10\baselineskip}

\subsubsection{Posterior bounds}\label{posterior-bounds}

\protect\phantomsection\label{result-proposition-a-1}{\textbf{Proposition
A.1 (posterior bounds; analytical).}} \emph{For arbitrary mixed
class-contingent orders on \([-1,1]\), the order-flow posterior
\(\mu_X\) and the posterior based on the price lie in
\([m(\pi),M(\pi)]\), where \(a_H\), \(a_L\), \(\mu_X\), \(m(\pi)\), and
\(M(\pi)\) are defined in \eqref{eq:paper-7}.}

For any feasible orders \(q,q'\) the triangle inequality gives
\(\big||x-q|-|x-q'|\big|\le|q-q'|\le2\), which is \eqref{eq:paper-8}.
Integrating against \(d\sigma_H(q)\,d\sigma_L(q')\) preserves both
inequalities, so \(e^{-2/b}a_L\le a_H\le e^{2/b}a_L\) for arbitrary
conditional mixtures. Since
\(\log[\mu_X/(1-\mu_X)]=\log[\pi/(1-\pi)]+\log(a_H/a_L)\), the log odds
of the posterior lie within \(2/b\) of the log odds of the prior, which
is \(m(\pi)\le\mu_X\le M(\pi)\). No monotone or pure strategy is
assumed. Because the price is a function of \(X\), the price posterior
is \(\mathbb E[\mu_X\mid P]\) by the tower property and inherits the
interval. The bounds satisfy \(1-M(\pi)=m(1-\pi)\), so
\(\kappa(\pi)=m(\min\{\pi,1-\pi\})\), and at \(\pi=1/2\) they are \(m\)
and \(M=1-m\). Online Appendix A.1 gives the measure-theoretic version,
including the null-set invariance under unilateral deviations.
\(\square\)

\Needspace{10\baselineskip}

\subsubsection{Price sufficiency and residual
profits}\label{price-sufficiency-and-residual-profits}

\protect\phantomsection\label{result-proposition-a-2}{\textbf{Proposition
A.2 (price sufficiency and residual profits; analytical).}}
\emph{Suppose \(\Delta_T>0\). In every candidate equilibrium the
observed price reveals \(\mu_X\) almost surely. Entry and competitive
pricing take the form \eqref{eq:paper-9}, and the residual advantages
of an informed buyer of shares in class \(H\) and of an informed short
seller in class \(L\) are given by \eqref{eq:paper-11}, with
\(\rho\kappa(\pi)\Delta_T\le A_H(x),A_L(x)\le\Delta_T\).}

\emph{Floor.} The low-cost challenger pays nothing, and its gross profit
at every belief is an expectation of nonnegative auction rents, so it
prepares at every price, by the tie rule when that profit is zero.
Cost-averaged entry conditional on the price, \(e(P)\), is therefore at
least \(\rho\). This holds before any division and needs no condition on
beliefs.

\emph{Gross profit at a price.} The challenger's value depends on flow
only through its class, so its expected gross profit given the price is
\(\sum_j\Pr(j\mid P)g_j=B(\Pr(H\mid P))\). Its preparation decision
therefore depends on the price only through its posterior.

\emph{Inversion.} Entry depends on \((P,C)\), and \(C\) is independent
of the class and of flow, so
\(\mathbb E[V_T\mid j,x]=t_0+e(P(x))(t_j-t_0)\). Competitive pricing
gives the first line of \eqref{eq:paper-10} pointwise in flow, with
\(e(P)\) a function of the price. Since
\(e(P)\Delta_T\ge\rho\Delta_T>0\), the second line of
\eqref{eq:paper-10} expresses \(\mu_X\) as a Borel function of \(P\).
Conditioning \(\mu_X\) on the price therefore leaves it unchanged, and
\(\Pr(H\mid P)=\mathbb E[\mu_X\mid P]=\mu_X\) almost surely. Entry is
then \eqref{eq:paper-9} with the buyer's own posterior.

\emph{Construction.} Put \(w_L=t_L-t_0\ge0\). Under the restrictions of
\hyperref[sec-values-preparation-and-the-sale]{Section 2.2},
\(g_H\ge g_L\), so entry is nondecreasing in the belief. For
\(\mu_2>\mu_1\), \eqref{eq:paper-9} gives

\[
P(\mu_2)-P(\mu_1)
=e(\mu_2)\Delta_T(\mu_2-\mu_1)+[e(\mu_2)-e(\mu_1)](w_L+\Delta_T\mu_1)
\ge\rho\Delta_T(\mu_2-\mu_1)>0,
\tag{A.3}\label{eq:paper-a-3}
\]

because entry is nondecreasing and the bracket is nonnegative. The price
is therefore strictly increasing in the posterior, including across the
jump where high-cost preparation becomes worthwhile, and its inverse on
its image is measurable. A price value inside the skipped interval is
generated by no flow. A flat posterior tail maps into an actual price
atom, which is conditioned on as an atom.

\emph{Residuals.} Conditional on class \(j\) and flow \(x\), expected
terminal value is \(t_0+e(P(x))(t_j-t_0)\), in which the no-entry
proceeds and the public entry effect both appear. Subtracting the
rational price \(t_0+e(P(x))[w_L+\Delta_T\mu_X(x)]\) in class \(H\), and
reversing the subtraction in class \(L\), the terms \(t_0\) and
\(e\,w_L\) cancel and leave \eqref{eq:paper-11}. No unanticipated
premium is attached. This step uses neither the inversion nor the sign
of \(\Delta_T\). The bounds follow from \(\rho\le e\le1\) and
\(m(\pi)\le\mu_X\le M(\pi)\): \(A_H\ge\rho[1-M(\pi)]\Delta_T\) and
\(A_L\ge\rho m(\pi)\Delta_T\). If instead \(\Delta_T=0\), both residuals
vanish, every nonzero order loses its trading cost, and the posterior
stays at the prior, which the constant price then reveals. Online
Appendix A.3 supplies the conditional-probability argument. \(\square\)
